\documentclass[11pt]{article}
\usepackage{amsmath,amssymb,amsthm}
\usepackage[margin=1in]{geometry}

\newtheorem{theorem}{Theorem}
\newtheorem{lemma}[theorem]{Lemma}

\title{Problem 433: Steps in Euclid's Algorithm}
\author{Project Euler Solutions}
\date{}

\begin{document}
\maketitle

\section{Problem Statement}
Let $E(x, y)$ be the number of division steps in the Euclidean algorithm: $E(x, 0) = 0$, $E(x, y) = 1 + E(y, x \bmod y)$ for $y > 0$. Define $S(N) = \sum_{1 \leq x, y \leq N} E(x, y)$. Find $S(5 \times 10^6)$.

\section{Mathematical Foundation}

\begin{theorem}[Termination Bound]
For all positive integers $x, y$, the Euclidean algorithm terminates in at most $\lfloor \log_\varphi(\sqrt{5}\cdot\min(x,y)) \rfloor$ steps, where $\varphi = (1+\sqrt{5})/2$.
\end{theorem}

\begin{proof}
By Lam\'e's theorem, if $E(x,y) = k$ with $x \geq y > 0$, then $y \geq F_{k+1}$ where $F_n$ is the $n$-th Fibonacci number. Since $F_n \geq \varphi^{n-2}/\sqrt{5}$, the bound follows.
\end{proof}

\begin{lemma}[Near-Symmetry]
For $x < y$, we have $E(x, y) = 1 + E(y, x)$, and $E(x, x) = 1$ for all $x \geq 1$.
\end{lemma}

\begin{proof}
If $x < y$, then $x \bmod y = x$, so $E(x, y) = 1 + E(y, x)$. For $x = y$: $E(x, x) = 1 + E(x, 0) = 1$.
\end{proof}

\begin{theorem}[Decomposition of $S(N)$]
\[
S(N) = 2\sum_{x=2}^{N}\sum_{y=1}^{x-1} E(x, y) + N.
\]
\end{theorem}

\begin{proof}
Partition the sum into diagonal ($x = y$), upper-triangular ($x > y$), and lower-triangular ($x < y$) parts. By Lemma~1, the lower-triangular contribution for each pair $(x, y)$ with $x < y$ is $E(x, y) = 1 + E(y, x)$. Combining the upper and lower parts and adding the diagonal (each contributing 1) yields the formula.
\end{proof}

\begin{theorem}[Continued Fraction Interpretation]
The step count $E(x, y)$ equals the length of the continued fraction expansion of $x/y$ (for $x > y > 0$ and $\gcd(x,y) = 1$). The sum $S(N)$ can be decomposed over $\gcd$ classes: pairs with $\gcd(x,y) = g$ contribute the same step count as the coprime pair $(x/g, y/g)$.
\end{theorem}

\begin{proof}
The Euclidean algorithm on $(x, y)$ produces quotients $q_1, q_2, \ldots, q_k$ which are precisely the partial quotients of the continued fraction $x/y = [q_1; q_2, \ldots, q_k]$. Since $\gcd$ is homogeneous, $E(gx', gy') = E(x', y')$.
\end{proof}

\section{Algorithm}
\begin{enumerate}
    \item Use the decomposition $S(N) = 2\sum_{x>y} E(x,y) + N$.
    \item For each pair $(x, y)$ with $1 \leq y < x \leq N$, compute $E(x, y)$ via the Euclidean algorithm.
    \item Exploit continued fraction structure and $\gcd$-class summation for sub-quadratic acceleration.
\end{enumerate}

\section{Complexity Analysis}
\begin{itemize}
    \item \textbf{Time:} $O(N^2 \log N)$ brute-force; sub-quadratic with continued fraction optimizations.
    \item \textbf{Space:} $O(1)$ auxiliary for brute-force; $O(N)$ for precomputation approaches.
\end{itemize}

\section{Answer}

\[
\boxed{326624372659664}
\]

\end{document}
